\documentclass[11pt]{article}

\usepackage[margin=1in]{geometry}
\usepackage{amsmath,amssymb,amsthm}
\usepackage{graphicx}
\usepackage[hidelinks]{hyperref}

\newtheorem{theorem}{Theorem}
\newtheorem{corollary}{Corollary}

\newcommand{\K}{\mathbb K}
\newcommand{\X}{\mathcal X}
\newcommand{\supp}{\operatorname{supp}}
\newcommand{\rank}{\operatorname{rank}}

\title{Exact Banach Support-Sum Criteria for Tao-Type Uncertainty Principles}
\author{Run \texttt{fa\_banach\_001}, agent \texttt{agent\_lane\_06}}
\date{June 28, 2026}

\begin{document}
\maketitle

\section*{Source and Open Problem}

The source paper is K. Mahesh Krishna, \emph{Functional Donoho-Stark-Elad-Bruckstein-Ricaud-Torresani Uncertainty Principle}, arXiv:2304.03324. In Section 3, after recalling Tao's prime-dimensional uncertainty principle for the Fourier transform on $\mathbb C^p$, the paper asks:
\[
\text{``What is the version Tao uncertainty principle for Banach spaces of prime dimension?''}
\]

\begin{center}
\includegraphics[width=0.94\linewidth]{figures/open_problem_crop.png}
\end{center}

\section*{Result}

For arbitrary finite p-Schauder frames, every support-sum lower bound is governed by a deletion/spanning condition on the remaining analysis functionals. In the basis case this becomes an especially simple full-spark criterion: the Tao support-sum bound is equivalent to the transition matrix being full spark. Thus the Banach-space replacement for the Fourier/primality mechanism is not prime dimension by itself, but the full-spark property of the two coordinate systems.

\section*{Setup}

Let $\X$ be a $d$-dimensional Banach space over $\K=\mathbb R$ or $\mathbb C$. Let
\[
(\{f_j\}_{j=1}^n,\{\tau_j\}_{j=1}^n),\qquad
(\{g_k\}_{k=1}^m,\{\omega_k\}_{k=1}^m)
\]
be p-Schauder frames in the sense of the source paper. Only their analysis functionals are needed below. The analysis maps
\[
\theta_f x=(f_j(x))_{j=1}^n,\qquad \theta_g x=(g_k(x))_{k=1}^m
\]
are injective; for p-Schauder frames they are isometries.

For $A\subseteq\{1,\dots,n\}$ and $B\subseteq\{1,\dots,m\}$ define
\[
E(A,B)=\bigcap_{j\notin A}\ker f_j\cap \bigcap_{k\notin B}\ker g_k.
\]
Equivalently, $E(A,B)$ is the common kernel of all analysis functionals that remain after deleting the coordinates in $A$ and $B$.

\section*{Proof Intuition}

For general p-Schauder frames, a vector has small support exactly when it lies in many coordinate kernels. Thus a lower bound on
\[
|\supp(\theta_f x)|+|\supp(\theta_g x)|
\]
is equivalent to saying that no nonzero vector can survive after all but a small number of coordinates are forced to vanish. In the basis case this deletion condition is the familiar full-spark condition on the transition matrix. Tao's finite cyclic prime theorem rests on the corresponding algebraic fact for the prime Fourier matrix: every square minor is nonzero.

\begin{theorem}[Exact frame support-sum criterion]
Let $L\ge 1$ be an integer. With the notation above, the following are equivalent.
\begin{enumerate}
\item For every $x\in \X\setminus\{0\}$,
\[
|\supp(\theta_f x)|+|\supp(\theta_g x)|\ge L.
\]
\item For every $A\subseteq\{1,\dots,n\}$ and $B\subseteq\{1,\dots,m\}$ with
\[
|A|+|B|<L,
\]
one has $E(A,B)=\{0\}$.
\item Equivalently, for every such $A,B$, the family
\[
\{f_j:j\notin A\}\cup \{g_k:k\notin B\}
\]
spans $\X^*$.
\end{enumerate}
\end{theorem}

\begin{proof}
For $x\ne0$, put
\[
A=\supp(\theta_f x),\qquad B=\supp(\theta_g x).
\]
Then $x\in E(A,B)$. Therefore if $|A|+|B|<L$ and $E(A,B)=\{0\}$ for all such pairs, no nonzero $x$ can have support sum below $L$.

Conversely, if $E(A,B)$ contains a nonzero vector $x$ for some pair with $|A|+|B|<L$, then all nonzero coordinates of $\theta_f x$ lie in $A$ and all nonzero coordinates of $\theta_g x$ lie in $B$. Hence
\[
|\supp(\theta_f x)|+|\supp(\theta_g x)|\le |A|+|B|<L,
\]
contradicting the support-sum bound.

Finally, a family of linear functionals spans $\X^*$ if and only if its common kernel is zero, giving the equivalence of (2) and (3).
\end{proof}

\begin{corollary}[Tao-type frame criterion]
For arbitrary p-Schauder frames on a $d$-dimensional Banach space, the support-sum bound
\[
|\supp(\theta_f x)|+|\supp(\theta_g x)|\ge d+1,\qquad x\ne0,
\]
holds if and only if, after deleting any total of at most $d$ coordinates from the two analysis families, the remaining functionals still span $\X^*$.
\end{corollary}

\section*{Basis Specialization}

Now suppose $n=m=d$ and the two frames are bases with biorthogonal coordinate functionals. Write
\[
A=(g_k(\tau_j))_{k,j=1}^d.
\]
Then $(g_k(x))=A(f_j(x))$. Call $A$ \emph{full spark} if every square minor of $A$ is nonzero.

\begin{theorem}[Basis full-spark criterion]
For two bases of $\X$, the following are equivalent.
\begin{enumerate}
\item For every $x\in \X\setminus\{0\}$,
\[
|\supp(f_j(x))|+|\supp(g_k(x))|\ge d+1.
\]
\item The transition matrix $A=(g_k(\tau_j))_{k,j=1}^d$ is full spark.
\end{enumerate}
\end{theorem}

\begin{proof}
This follows from the preceding corollary, but we spell out the standard matrix form. If $A$ is not full spark, choose $R,S\subseteq\{1,\dots,d\}$ with $|R|=|S|=s$ such that $A_{R,S}$ is singular. There is a nonzero coordinate vector $c$ supported in $S$ with $A_{R,S}c_S=0$. For $x=\sum_j c_j\tau_j$ one has
\[
|\supp(f_j(x))|\le s,\qquad |\supp(g_k(x))|\le d-s,
\]
so the support sum is at most $d$.

Conversely, suppose there is a nonzero $x$ with support sets
\[
S=\supp(f_j(x)),\qquad T=\supp(g_k(x)).
\]
If $s+t=|S|+|T|\le d$, then every row in $R_0=\{1,\dots,d\}\setminus T$ annihilates the nonzero vector $c_S=(f_j(x))_{j\in S}$ through $A_{R_0,S}$. Since $|R_0|=d-t\ge s=|S|$, choose $R\subseteq R_0$ with $|R|=s$. Then
\[
A_{R,S}c_S=0
\]
with $c_S\ne0$, so the square minor $A_{R,S}$ is singular. Therefore $A$ is not full spark.
\end{proof}

\begin{corollary}[Prime-dimensional Banach formulation]
Let $r>1$ be prime. For two p-orthonormal bases in an $r$-dimensional Banach space, the Tao-type bound
\[
|\supp(\theta_f x)|+|\supp(\theta_g x)|\ge r+1,\qquad x\ne0,
\]
holds if and only if their transition matrix $(g_k(\tau_j))$ is full spark.
\end{corollary}

\begin{proof}
Apply the theorem with $d=r$. The proof does not use primality; in the classical Fourier case, primality enters through Tao's theorem that the prime Fourier matrix is full spark.
\end{proof}

\begin{corollary}[Prime dimension alone is insufficient]
There is no dimension-only Tao principle for arbitrary Banach p-orthonormal bases. For every prime $r>1$ and every $1<p<\infty$, take $\X=\ell_p^r$ and take the same canonical p-orthonormal basis twice. For $x=e_1$,
\[
|\supp(\theta_f x)|+|\supp(\theta_g x)|=1+1=2<r+1.
\]
\end{corollary}

\section*{Scope and Limitations}

This is a strong partial answer to the source problem because Problem 3.2 asks broadly for a ``version'' of Tao's theorem, not for a single formal conjecture. The packet settles the natural support-sum formulation for arbitrary finite p-Schauder frames and gives the clean full-spark specialization for p-orthonormal bases. It also shows that the prime-dimensional hypothesis alone cannot replace the Fourier full-spark mechanism.

The result does not claim to choose a unique canonical Banach-space analogue among possible entropy, product, weighted-support, or recovery-theoretic Tao variants. It gives the exact answer for the support-sum interpretation suggested by the source's use of support cardinalities.

\section*{Novelty and Literature Check}

The run indexes were searched for \texttt{2304.03324}, the source title, Tao uncertainty terminology, prime dimension, and full spark. No duplicate packet was found. Web searches on June 28, 2026 for the exact Problem 3.2 wording, ``Functional Tao uncertainty principle'', and ``Tao uncertainty principle Banach spaces prime dimension'' surfaced the source paper and adjacent Krishna followups but no later explicit answer to Problem 3.2.

\section*{Human Review Recommendation}

Review as a likely valid strong partial result. The key verification is the elementary equivalence between support-sum lower bounds and the deletion/spanning condition. The full-spark basis corollary and the warning that prime dimension alone is insufficient then follow immediately.

\begin{thebibliography}{9}

\bibitem{Krishna}
K. Mahesh Krishna,
\emph{Functional Donoho-Stark-Elad-Bruckstein-Ricaud-Torresani Uncertainty Principle},
arXiv:2304.03324, 2023.

\bibitem{Tao}
Terence Tao,
\emph{An uncertainty principle for cyclic groups of prime order},
Math. Res. Lett. 12 (2005), no. 1, 121--127.

\end{thebibliography}

\end{document}
